\documentclass[instructions.tex]{subfiles}
\begin{document}
 
\chapter{Pour être pénitent, il faut vivre dans le gémissement et le regret.}
 
\primo Il sert peu d’affliger le corps, si l’esprit et le cœur ne sont pas affligés. La tristesse du Sauveur à la vue des péchés du monde, sa douleur profonde au jardin des Oliviers, est une puissante leçon. Un Dieu qui pleure mes péchés, qui s’afflige sur l’état de mon âme, me fait comprendre que je ne puis trop affliger mon cœur, ni répandre trop de larmes.
 
Les gémissemens ont fait la principale occupation des saints sur la terre. David, quoique assuré du pardon, gémissait sans cesse ; le souvenir d’avoir perdu son Dieu lui faisait répandre nuit et jour des torrens de larmes. Saint Pierre ne pécha qu’une fois, et il en pleura toujours\footnote{Semel negavit, et semper flevit ; sæpe negamus, et nunquam flemus. S. Aug.} ; et nous, au contraire, nous péchons souvent, et nous ne pleurons jamais. En quelque lieu que je sois, disait un solitaire, je ne vois que mes péchés, je me regarde comme une victime de l’enfer, où je vois une infinité d’âmes moins coupables que moi. Alors je me jette la face contre terre, je soupire, je pleure devant mon juge.
 
\secundo Voulez-vous obtenir la miséricorde de Dieu ? que vos péchés soient toujours présens à votre esprit\footnote{Peccatum meum contrà me est semper. P. 50. 4.} . Dites sans cesse avec un prophète : Ah ! Seigneur, qu’ai-je fait\footnote{Quid feci ? 1. Reg. 17. 29.} ? Dites avec saint Augustin : Ah ! mon Dieu, beauté toujours ancienne et toujours nouvelle, c’est bien tard que je vous ai connu, c’est bien tard que je vous ai aimé. Dites avec l’enfant prodigue : Mon père, j’ai péché contre le ciel et contre vous, je ne mérite plus d’être appelé votre enfant.
 
Sainte Thaïs, fameuse pécheresse, dit à saint Paphnuce, que pendant trois ans elle n’avait osé proférer le nom de Dieu, tant elle se sentait criminelle ; et qu’elle s’était contentée de dire sans cesse en gémissant : O vous qui m’avez formée, ayez pitié de moi ! Allez, lui dit le saint homme, vos gémissemens, plus que vos autres pénitences, ont engagé Dieu à vous faire miséricorde.
 
Ne croyez jamais avoir assez pleuré vos péchés, dit saint Augustin ; que la douleur et la confusion ne finissent qu’avec votre vie. Un seul péché, dit Tertullien, mérite d’être pleuré éternellement\footnote{Semel peccasse, satis est ad fletus æternos. Tertul.} . Lorsqu’une âme a sans cesse ses fautes devant les yeux, et qu’elle en gémit, elle peut s’assurer que Dieu lui fera miséricorde. Rien ne touche plus Dieu, que le repentir et le sacrifice d’un cœur affligé\footnote{Sacrificium Deo, spiritus contribulatus. Ps. 50. 18.} .
 
\section{Résolutions}
 
\primo Et moi aussi, à l’exemple de tant des saints pénitens, je me ressouviendrai souvent de mes péchés passés, pour les détester et m’efforcer de les laver dans les larmes d’un repentir amer.
 
\secundo Et je crierai fréquemment miséricorde vers celui que j’ai tant offensé. 

\prayer{O mon Dieu ! soyez propice à ce grand pécheur.}
 
\end{document}
